% Delta from §10.27. Insert at the end of 05_ciclo_catalan.tex.
% The affine update and balance remain in their original source sections.
\subsection{Return traces and moments of the constitutive quotient}
\label{amp-afin-sec}

The register update and the cycle response share the dodecaphasic
transport and preserve different kinds of information.
The signed event is obtained from the direction and orientation of each
window of the APP--TRIT--TPK history. Its vector increment and
balance~\eqref{mem:balance} retain the interaction with the previous
register. The construction below extends the spectral reading of the
same cycle, once its incidence subspaces have been fixed. This order
allows return, determinant and moments to be composed while preserving
the history that precedes them.

Retain the operator $C$ and projectors of
\eqref{iii:eq:projectors34}, and define
\begin{equation}
 E_{\rm cyc}=P_3^{\mathrm{cyc}}-P_4^{\mathrm{cyc}},\qquad
 b_n=\operatorname{Tr}(C^n E_{\rm cyc}),\qquad n\geq0.
 \label{amp-afin-trazas}
\end{equation}
The symbol $b_n$ distinguishes this spectral trace from the signed
event $a_m$ of a window. The trace uses the transport operator; the
event also uses the trajectory actually traversed.

\begin{proposition}[Trace sequence produced by incidence]
\label{amp-afin-prop-tabla}
One has
\begin{equation}
 \begin{gathered}
 b_n=3\mathbf1_{3\mid n}-4\mathbf1_{4\mid n},\\
 (b_1,\ldots,b_{12})=(0,0,3,-4,0,3,0,-4,3,0,0,-1).
 \end{gathered}
 \label{amp-afin-tabla}
\end{equation}
The least period is twelve and the sum over a period is zero. For
$B_N=\sum_{n=1}^{N}b_n$,
\begin{equation}
 B_N=3\lfloor N/3\rfloor-4\lfloor N/4\rfloor,\qquad |B_N|\leq3.
 \label{amp-afin-parciales}
\end{equation}
\end{proposition}
\begin{proof}
On $H_d$, the operator $C$ cyclically permutes the $d$ coordinate
classes. Its $n$th power has $d$ fixed points if $d$ divides $n$,
and none otherwise. Since $P_d^{\mathrm{cyc}}$ commutes with $C$,
$\operatorname{Tr}(C^nP_d^{\mathrm{cyc}})$ is the trace of this
restriction. Subtracting the cases $d=3,4$ proves the formula.
The least positive period is twelve: $b_n=-1$ occurs exactly at
multiples of twelve. Over one period the contributions are
$3\cdot4-4\cdot3=0$. The finite sum counts multiples and gives
\eqref{amp-afin-parciales}. Its values for $N=0,\ldots,11$ are
$(0,0,0,3,-1,-1,2,2,-2,1,1,1)$ and repeat after adding twelve.
\end{proof}

\begin{theorem}[Logarithmic reconstruction of the constitutive reader]
\label{amp-afin-thm-logdet}
For $|s|<1$, the quotient $f$ of~\eqref{iii:eq:trace-det} satisfies
\begin{align}
 -\log f(s)&=\sum_{n\geq1}\frac{b_n}{n}s^n,\label{amp-afin-logdet}\\
 -\frac{s f'(s)}{f(s)}
 &=\frac{3s^3}{1-s^3}-\frac{4s^4}{1-s^4}.
 \label{amp-afin-logder}
\end{align}
The holomorphic branch is fixed by $\log f(0)=0$.
For $0<s<1$, it agrees with the real logarithm.
\end{theorem}
\begin{proof}
The polynomials $1-s^3$ and $1-s^4$ do not vanish in the open disc.
The expansion $-\log(1-w)=\sum_{k\geq1}w^k/k$, obtained by
integrating the geometric series from zero, converges uniformly on
compact discs. Substitution of $w=s^3,s^4$ and subtraction give
$-\log f(s)=\sum_{k\geq1}s^{3k}/k-\sum_{k\geq1}s^{4k}/k$.
The coefficient of $s^n$ is $b_n/n$.
Local uniform convergence permits differentiation and summation of
the two geometric series, proving~\eqref{amp-afin-logder}.
The condition at zero fixes the integration constant and the branch.
\end{proof}

Thus $b_0/12=-1/12$ records the normalized dimension defect,
whereas the full sequence $b_n$ reconstructs a function.
These are readings of different scope on the same pair of subspaces.
Upon substituting $s=q_\pm^{30}$, the arguments remain outputs of
the preceding angular construction; the expansion does not select them.

\begin{theorem}[Depth moments and Mellin representation]
\label{amp-afin-thm-momentos}
The series
\begin{equation}
 M_b(z)=\sum_{n\geq1}\frac{b_n}{n^z}
 \label{amp-afin-momento}
\end{equation}
converges locally uniformly for $\operatorname{Re}z>0$.
Writing $\sigma=\operatorname{Re}z>0$, its remainder satisfies
\begin{equation}
 \left|M_b(z)-\sum_{n=1}^{N}\frac{b_n}{n^z}\right|
 \leq3\left(1+\frac{|z|}{\sigma}\right)(N+1)^{-\sigma}.
 \label{amp-afin-cota}
\end{equation}
In the domain of absolute convergence $\operatorname{Re}z>1$,
\begin{equation}
 M_b(z)=(3^{1-z}-4^{1-z})\zeta(z),\qquad
 \zeta(z)=\sum_{n\geq1}n^{-z}.
 \label{amp-afin-reconocimiento}
\end{equation}
For $\operatorname{Re}z>0$ one also has
\begin{equation}
 \Gamma(z)M_b(z)=
 \int_0^\infty t^{z-1}
 \left(\frac3{e^{3t}-1}-\frac4{e^{4t}-1}\right)dt,
 \label{amp-afin-mellin}
\end{equation}
where $\Gamma(z)=\int_0^\infty u^{z-1}e^{-u}\,du$.
\end{theorem}
\begin{proof}
Summation by parts and $|B_n|\leq3$ give, after taking the upper limit,
\[
 \sum_{n>N}b_n n^{-z}
 =-B_N(N+1)^{-z}
   +\sum_{n>N}B_n\bigl(n^{-z}-(n+1)^{-z}\bigr).
\]
The modulus of the difference of powers is bounded by
$|z|\int_n^{n+1}x^{-\sigma-1}dx$. This proves
\eqref{amp-afin-cota}, convergence and uniformity on compact subsets.
For $\sigma>1$, absolute convergence permits multiples of three
and four to be grouped separately, giving
\eqref{amp-afin-reconocimiento}.

The geometric series of the transport, evaluated at $s=e^{-t}$,
produces
\[
 \sum_{n\geq1}b_n e^{-nt}
 =\frac3{e^{3t}-1}-\frac4{e^{4t}-1}=:K_b(t).
\]
For $\sigma>1$, termwise integration is justified by
$\sum|b_n|n^{-\sigma}<\infty$. Substitution of $u=nt$ gives
$\Gamma(z)n^{-z}$ and proves~\eqref{amp-afin-mellin} on that
half-plane. At the origin the $1/t$ terms cancel:
$K_b(t)=1/2+O(t)$. At infinity, $K_b(t)=O(e^{-3t})$.
The integral is therefore holomorphic for $\operatorname{Re}z>0$
by local uniform convergence. The series is also holomorphic there
by the preceding bound; the identity theorem extends equality
to that domain.
\end{proof}

For $0<\operatorname{Re}z\leq1$, increasing depth fixes the
convergent ordering of the series. Its realization as the absolute
trace of the diagonal operator applies for $\operatorname{Re}z>1$.

\begin{corollary}[Evaluations compatible with the constitutive endpoint]
\label{amp-afin-cor-extremos}
One has
\begin{equation}
 M_b(2)=\frac{\zeta(2)}{12},\qquad
 M_b(1)=\log(4/3)=-\log f(1),\qquad f(1)=\frac34.
 \label{amp-afin-extremos}
\end{equation}
\end{corollary}
\begin{proof}
The first equality follows from $3^{-1}-4^{-1}=1/12$.
For the second, the exact partial sums are
$H_{\lfloor N/3\rfloor}-H_{\lfloor N/4\rfloor}$, with
$H_j=\sum_{k=1}^j1/k$ and $H_0=0$.
Integral comparison for $1/x$ between these two indices gives
the limit $\log(4/3)$. Cancelling $1-s$ before evaluating the
determinantal quotient gives $f(1)=3/4$.
\end{proof}

The affine register retains events and their interaction with previous
memory; the trace sequence retains the incidence of the two subspaces
at every power of transport; the moment applies a depth weighting
to that already-produced sequence. Its spectral periodicity coexists
with an event history that continues after phase return.
Replacing $C$ by its inverse preserves $b_n$, since $d\mid n$
is invariant under $n\mapsto-n$; the oriented resolvent coefficient
in~\eqref{iii:eq:resolvent}, with fixed marks, changes sign.
The joint reading preserves a distinction not recorded by the
scalar determinant alone. The quadratic form and coordinate sum
of the register retain their mathematical types; this composition
assigns them neither energy units nor electric-charge units.

\subsection{Residual composition and products of moments}
\label{amp-afin-sec-productos}

The subsequent composition of readers also retains the operation
performed on depth. In the realization
$\mathcal H=\ell^2(\mathbb N_{>0})$, set
$N\delta_n=n\delta_n$. Let $a,c$ be two already-produced periodic
residual tables and let $Q_a\delta_n=a(n)\delta_n$,
$Q_c\delta_n=c(n)\delta_n$. After extension to a common period,
the two readers act on the same index. For
$\operatorname{Re}z>1$, write
$M_a(z)=\operatorname{Tr}(N^{-z}Q_a)$ and similarly for $M_c$.

\begin{proposition}[Diagonal composition and tensor composition]
\label{amp-afin-prop-productos}
The following identities hold:
\begin{align}
 \operatorname{Tr}(N^{-z}Q_aQ_c)
 &=\sum_{n\geq1}\frac{a(n)c(n)}{n^z},
 \label{amp-afin-diagonal}\\
 M_a(z)M_c(z)
 &=\operatorname{Tr}_{\mathcal H\otimes\mathcal H}
 \bigl((N\otimes N)^{-z}(Q_a\otimes Q_c)\bigr)\notag\\
 &=\sum_{k\geq1}\frac{(a*c)(k)}{k^z},
 \label{amp-afin-tensor}\\
 (a*c)(k)&=\sum_{d\mid k}a(d)c(k/d),
 \qquad \operatorname{Re}z>1.
 \label{amp-afin-convolucion}
\end{align}
Every operator appearing inside these traces is trace class.
\end{proposition}
\begin{proof}
The diagonal product acts on $\delta_n$ by $a(n)c(n)$,
proving~\eqref{amp-afin-diagonal}.
On $\delta_m\otimes\delta_n$, the operator $N\otimes N$
has eigenvalue $mn$, and $Q_a\otimes Q_c$ acts by $a(m)c(n)$.
If $\sigma=\operatorname{Re}z>1$, the sum of the moduli of these
spectral coefficients is bounded by
$\|a\|_\infty\|c\|_\infty(\sum n^{-\sigma})^2<\infty$.
This bound establishes the trace-class condition and permits the
double sum to be factorized or grouped by $k=mn$.
Pairs with product $k$ are exactly $(d,k/d)$ with $d\mid k$.
The analogous bound with a single sum proves the diagonal case.
\end{proof}

For the readers in this section, taking $a(n)=b_n$ and
$c(n)=\langle v,C^nu\rangle$ gives, at $n=3$,
$a(3)c(3)=-3$, whereas $(a*c)(3)=3$.
This coefficient difference exhibits the information retained by
each operation. Diagonal composition retains a common depth;
tensor composition retains two depths before grouping by their product.
The tensor realization does not ascribe statistical independence
to the HMT histories that produced the readers.
Regularized derivatives and normalization changes of other moments
are subsequent operations distinct from the two compositions proved here.
